% ---------------------------------------------------------------- Introduction (P1)
\section{Introduction}

Pricing an American option means deciding when to exercise it, which makes it an
\emph{optimal stopping} problem. The classical numerical methods (binomial trees,
the least-squares Monte Carlo method of Longstaff and Schwartz~\cite{ls2001} and Deep
Optimal Stopping~\cite{becker2019}) all solve it by backward induction. They assume that
the price process is Markovian, so that today's state is enough to decide.

Real markets are not Markovian in the price alone: volatility is persistent and,
empirically, \emph{rough}~\cite{bfg2016}. When the past matters, backward induction needs
an ever larger state, or it carries approximation errors from one date into the
previous one. Damera Venkata and Bhattacharyya~\cite{ospg2023} proposed the
\emph{Optimal Stopping Policy Gradient} (OSPG), which trains a recurrent network to
choose stopping times directly and needs no backward recursion.

We (i) replicate OSPG on the paper's Markovian option benchmarks, (ii) price American
puts under the rough Bergomi model, where the price is non-Markovian, and (iii) backtest
learned exit-timing rules on market data, comparing each with the classical methods.
